%% KTH Document Template — Example
%% Based on kth-document.cls, implementing the KTH Graphical Profile (Sept. 2023)
%%
%% ---- Compilation -------------------------------------------------------
%%   pdflatex example.tex      standard, uses TeX Gyre fonts as substitutes
%%   xelatex  example.tex      exact KTH fonts (Figtree + Georgia) if installed
%%
%% ---- Logo --------------------------------------------------------------
%%   The logo is shown only on the first page (in the title block).
%%   The class looks for the file set by \kthlogopath in this order:
%%     .pdf  →  .png  →  .eps
%%   Default filename: KTH_logo_RGB_bla  (the standard KTH blue logo PNG)
%%   To use a different file:
%%     \renewcommand{\kthlogopath}{/path/to/your-logo}   (no extension)
%%   To adjust logo size:
%%     \renewcommand{\kthlogoheight}{1.8cm}
%%
%% ---- Class options -----------------------------------------------------
%%   english     — English locale: babel, standard decimal point  [default]
%%   swedish     — Swedish locale: babel, decimal comma via icomma
%%   titlepage   — full separate title page instead of inline title block
%%   12pt, etc.  — any article class option is passed through
%%
%% ---- Metadata commands -------------------------------------------------
%%   \title{...}        document title                           (required)
%%   \subtitle{...}     subtitle line below the title           (optional)
%%   \doctype{...}      document type label, e.g. "Activity Report"
%%   \author{...}       author name(s)
%%   \affiliation{...}  department / school, shown next to author
%%   \date{...}         date; use \today or a fixed string
%%   \version{...}      version string, shown as "(vX.Y)" after date
%%   \shorttitle{...}   short title shown in header on pages 2+
%%
%% ---- KTH colors --------------------------------------------------------
%%   All colors from the KTH Grafisk manual (v1.2, 2024) are defined and
%%   can be used anywhere with \textcolor{name}{...} or \colorbox{name}{...}.
%%
%%   Primary palette:
%%     kthwhite        #FFFFFF
%%     kthsand         #EBE5E0   warm grey
%%     kthblue         #004791   KTH brand blue — headings, rules, links
%%     kthskyblue      #6298D2
%%     kthnavy         #000061   navy — logo wreath color
%%     kthlightblue    #DEF0FF   tint (used for kthbox background)
%%     kthdigitalblue  #0029ED   screen only; use kthblue for print
%%
%%   Functional palette — for diagrams, charts, reports.
%%   Each color family has three variants: dark / mid / light.
%%   Naming: kth{dark,<empty>,light}{family}
%%
%%     family    dark (#)    mid (#)     light (#)
%%     green     0D4A21      4DA060      C7EBBA
%%     teal      1C434C      339C9C      B2E0E0
%%     brick     78001A      E86A58      FFCCC4
%%     yellow    A65900      FFBE00      FFF0B0
%%     gray      323232      A5A5A5      E6E6E6   (grey spellings also work)
%%
%%   Web-only text colors: kthbrokenblack (#212121), kthbrokenwhite (#FCFCFC)
%%   Aliases: kthaccent = kthblue,  kthmuted = kthsand
%%
%% ---- Callout boxes -----------------------------------------------------
%%   \begin{kthbox}...\end{kthbox}
%%       Light-blue highlighted box — use for summaries, key points, notices.
%%   \begin{kthnotebox}...\end{kthnotebox}
%%       Sand-colored box — use for notes, caveats, supplementary info.
%%
%% ---- Inline helpers ----------------------------------------------------
%%   \kthhl{text}      bold + KTH blue — highlight a keyword inline
%%
%% ---- Adding packages ---------------------------------------------------
%%   Add \usepackage{...} in the preamble below as usual.
%%   The class loads: geometry, xcolor, graphicx, hyperref, fancyhdr,
%%   titlesec, parskip, mdframed, babel, csquotes, iftex, microtype (pdftex).
%%   Safe to add: amsmath, amssymb, booktabs, biblatex, siunitx, tikz, etc.

\documentclass[english,11pt]{kth-document}

% ---- Optional: adjust logo path ----
% \renewcommand{\kthlogopath}{/path/to/kth-logo}

% ---- Metadata ----
\title{Project Title}
\subtitle{A short, informative subtitle}
\doctype{Activity Report}          % e.g. Project Description, Report, Memo
\author{Firstname Lastname}
\affiliation{School of Engineering Sciences (SCI), KTH}
\date{\today}
\version{1.0}                      % optional; omit to hide version
\shorttitle{Project Title}         % appears in header on pages 2+

% ---- Add packages as needed ----
% \usepackage{amsmath, amssymb}    % math
% \usepackage{booktabs}            % nice tables
% \usepackage{biblatex}            % bibliography

% Colored swatch box used in the color reference section below
\newcommand{\swatch}[1]{%
  {\setlength{\fboxsep}{0pt}\setlength{\fboxrule}{0.4pt}%
   \fbox{\colorbox{#1}{\rule{0pt}{0.9em}\hspace{2.2em}}}}%
}

\begin{document}

\maketitle

% ---- Optional abstract / summary box ----
\begin{kthbox}
  \textbf{Summary.}
  A brief summary of the document goes here. This highlighted box uses the
  KTH light blue and is useful for abstracts, key findings, or important notices. 
\end{kthbox}


\section{Background}

This document is an example of the \textbf{kth-document} LaTeX class, which
implements the KTH Graphical Profile (updated September 2023). It provides a
clean, professional look suitable for activity reports, project descriptions,
internal memos, and similar documents.

The class is built on the standard \texttt{article} class and stays out of the
way of additional packages you may need.


\section{Objectives}

\subsection{Primary goals}

Write a short description of the main objectives. Use bullet lists where
appropriate:

\begin{itemize}
  \item First objective — concise and specific.
  \item Second objective — with measurable outcomes.
  \item Third objective — connecting to broader context.
\end{itemize}

\subsection{Secondary goals}

Additional or supporting goals can go here. The \kthhl{highlighted keyword}
command draws attention to key terms without being distracting.


\section{Activities and Timeline}

\begin{center}
\begin{tabular}{llp{8cm}}
  \hline
  \textbf{Phase} & \textbf{Period} & \textbf{Description} \\
  \hline
  Planning   & Jan–Feb 2026 & Requirements gathering and stakeholder alignment. \\
  Execution  & Mar–May 2026 & Core implementation and iterative testing. \\
  Evaluation & Jun 2026     & Assessment against success criteria. \\
  \hline
\end{tabular}
\end{center}


\section{Results}

Describe outcomes, findings, or deliverables here.

\begin{kthnotebox}
  \textbf{Note.}
  The sand-colored box (\texttt{kthnotebox}) is useful for notes,
  caveats, or supplementary information.
\end{kthnotebox}


\section{Conclusions}

Summarize key takeaways and any next steps or recommendations.


\section{KTH Color Palette}

All colors from the KTH Grafisk manual (v1.2, 2024) are available as named
colors. Use them with \verb|\textcolor{name}{...}|, \verb|\colorbox{name}{...}|,
or anywhere \textsf{xcolor} accepts a color name.

\medskip
\textbf{Primary palette}

\smallskip
\begin{tabular}{@{}llll@{}}
  \swatch{kthblue}        & \texttt{kthblue}        & \texttt{\#004791} & KTH brand blue — headings, rules, links \\[4pt]
  \swatch{kthskyblue}     & \texttt{kthskyblue}     & \texttt{\#6298D2} & Sky blue \\[4pt]
  \swatch{kthnavy}        & \texttt{kthnavy}        & \texttt{\#000061} & Navy — logo wreath \\[4pt]
  \swatch{kthlightblue}   & \texttt{kthlightblue}   & \texttt{\#DEF0FF} & Light blue tint (\texttt{kthbox} background) \\[4pt]
  \swatch{kthsand}        & \texttt{kthsand}        & \texttt{\#EBE5E0} & Sand / warm grey (\texttt{kthnotebox} background) \\[4pt]
  \swatch{kthwhite}       & \texttt{kthwhite}       & \texttt{\#FFFFFF} & White \\[4pt]
  \swatch{kthdigitalblue} & \texttt{kthdigitalblue} & \texttt{\#0029ED} & Digital blue — \emph{screen only, not for print} \\
\end{tabular}

\medskip
\textbf{Functional palette} — for diagrams, charts, and reports.
Each family has dark, mid, and light variants named
\texttt{kth\{dark,\textnormal{<empty>},light\}\{family\}},
e.g.\ \texttt{kthdarkgreen}, \texttt{kthgreen}, \texttt{kthlightgreen}.
Grey spellings (\texttt{kth{*}grey}) are also accepted.

\smallskip
\begin{tabular}{@{}lcccp{8cm}@{}}
  \textit{Family} & \textit{dark} & \textit{mid} & \textit{light} & \\[4pt]
  green  & \swatch{kthdarkgreen}  & \swatch{kthgreen}  & \swatch{kthlightgreen}
         & \texttt{\#0D4A21 / \#4DA060 / \#C7EBBA} \\[4pt]
  teal   & \swatch{kthdarkteal}   & \swatch{kthteal}   & \swatch{kthlightteal}
         & \texttt{\#1C434C / \#339C9C / \#B2E0E0} \\[4pt]
  brick  & \swatch{kthdarkbrick}  & \swatch{kthbrick}  & \swatch{kthlightbrick}
         & \texttt{\#78001A / \#E86A58 / \#FFCCC4} \\[4pt]
  yellow & \swatch{kthdarkyellow} & \swatch{kthyellow} & \swatch{kthlightyellow}
         & \texttt{\#A65900 / \#FFBE00 / \#FFF0B0} \\[4pt]
  gray   & \swatch{kthdarkgray}   & \swatch{kthgray}   & \swatch{kthlightgray}
         & \texttt{\#323232 / \#A5A5A5 / \#E6E6E6} \\
\end{tabular}

\begin{kthnotebox}
  \textbf{Web-only text colors.}
  \texttt{kthbrokenblack} (\texttt{\#212121}) and \texttt{kthbrokenwhite}
  (\texttt{\#FCFCFC}) are intended for body text on screen to improve
  readability. Do not use them for decorative elements or print.
\end{kthnotebox}


\section{Contact}

For questions about this document, contact
\href{mailto:name@kth.se}{name@kth.se}.

\end{document}
